%% Manuscript for the Special Issue "Statistics for Data Science and AI",
%% International Journal of Data Science and Analytics (Springer Nature).
%% Compile with: latexmk -pdf main.tex   (see Makefile)

\documentclass[pdflatex,sn-mathphys-num]{sn-jnl}

\usepackage{graphicx}
\usepackage{multirow}
\usepackage{amsmath,amssymb,amsfonts}
\usepackage{amsthm}
\usepackage{mathrsfs}
\usepackage[title]{appendix}
\usepackage{xcolor}
\usepackage{textcomp}
\usepackage{manyfoot}
\usepackage{booktabs}
\usepackage{algorithm}
\usepackage{algorithmicx}
\usepackage{algpseudocode}

\theoremstyle{thmstyleone}
\newtheorem{proposition}{Proposition}
\theoremstyle{thmstylethree}
\newtheorem{remark}{Remark}
\newtheorem{definition}{Definition}

\newcommand{\E}{\mathbb{E}}
\newcommand{\Var}{\operatorname{Var}}
\newcommand{\Cov}{\operatorname{Cov}}
\newcommand{\dbar}{\bar d}
\makeatletter
\newcommand{\tabinput}[1]{\@@input #1 }% primitive input, usable inside tabular
\makeatother

\raggedbottom
\gdef\orcidlogo{\includegraphics[height=1.6ex]{Orcidlogo.eps}}% scaled ORCID icon

\begin{document}

\title[Are Machine Learning Improvements Statistically Real?]{Are Machine Learning Improvements Statistically Real? Dependence-Aware Inference for Reliable Model Comparison on Spatiotemporal Data}

% --------------------------------------------------------------------------
% Authors
% --------------------------------------------------------------------------

\author*[1,2]{\fnm{Isaac} \sur{Touza}\,\orcid{https://orcid.org/0009-0003-6288-9458}}

\email{isaac\_touza@outlook.fr}

% --------------------------------------------------------------------------
% Affiliations
% --------------------------------------------------------------------------

\affil*[1]{\orgdiv{Department of Mathematics--Computer Science, Faculty of Sciences},
	\orgname{University of Maroua},
	\orgaddress{\city{Maroua}, \postcode{814}, \country{Cameroon}}}

\affil[2]{\orgdiv{Laboratoire de Recherche en Informatique (LaRI)},
	\orgname{University of Maroua},
	\orgaddress{\street{P.O. Box 46}, \city{Maroua}, \country{Cameroon}}}

\abstract{Claims that one machine learning model outperforms another usually rest on cross-validated scores and a test that assumes independent observations. On spatial and temporal data this assumption fails twice: random cross-validation leaks information between neighbouring observations, and the prediction errors of the compared models are themselves dependent. We study the second, less examined problem. We cast model comparison as inference on the mean of a dependent field of out-of-fold loss differentials with correlated fold-level training effects, and show that the paired $t$-test on fold scores underestimates the variance by exactly the average covariance between fold means. We propose DA-HAC, a dependence-aware test computed from a single blocked cross-validation run. It combines a product-kernel variance estimator with bandwidths calibrated on the empirical variogram of the loss differentials, a method-of-moments estimate of the training variance, and a Satterthwaite reference distribution. In simulations with exact control of the truth, including XGBoost learners whose expected performances are equal by construction, conventional procedures reject a true null hypothesis in up to $80\%$ of replicates, and random cross-validation halves the estimated gain of a genuinely better model under strong dependence. Blocked cross-validation with DA-HAC keeps the error rate near the nominal level under weak dependence and reduces it substantially otherwise, while diagnosing, through an effective sample size and a range ratio, the cases in which no reliable comparison is possible. On hourly weather, pedestrian counts and yield-monitor data, the conclusions of four of six comparisons depend on these choices. An open-source Python implementation is provided.}

\keywords{Model comparison, Cross-validation, Spatio-temporal dependence, HAC variance estimation, Variogram, Statistical significance, Trustworthy machine learning}

\maketitle

\section{Introduction}\label{sec:intro}

Empirical machine learning advances one comparison at a time. A new gradient-boosting configuration is reported to beat a random forest, a deep architecture to beat gradient boosting, one feature set to beat another. Such claims rest on a small number of summary statistics (an average accuracy, an F1-score, a root mean squared error) computed by cross-validation, and, in the better cases, on a statistical test or a confidence interval for the difference. The test is meant to answer a simple question: \emph{is the observed improvement larger than what the randomness of the data and of the validation procedure could produce on its own?}

The standard tools for that question~\cite{dietterich1998,nadeau2003,demsar2006,benavoli2017} were developed for observations that are independent and identically distributed (i.i.d.). Many of the data sets to which machine learning is now applied are not. Meteorological, environmental, agricultural, epidemiological and mobility data are collected at locations and times that are close to one another, and nearby observations share unmeasured drivers. Two consequences are well documented. First, random $K$-fold cross-validation places near-duplicates of test observations in the training set, so that it measures interpolation rather than prediction at new locations or times, and it is optimistic when the latter is the target~\cite{roberts2017,ploton2020,meyer2018,bergmeir2012}. Second, the \emph{errors} of a fitted model are themselves dependent, so that $n$ test observations carry much less information than $n$ independent ones; this second problem has been addressed for fixed forecasts and spatial predictions~\cite{diebold1995,hering2011}, but rarely for learning procedures evaluated by cross-validation.

The first problem has a well-known remedy: blocked (spatial, temporal or spatiotemporal) cross-validation~\cite{roberts2017,valavi2019,linnenbrink2024}. The second problem is the subject of this paper. Blocking changes \emph{what} is estimated; it does not tell us \emph{how precisely} it is estimated. A practitioner who computes a blocked cross-validation score for two models and then applies a paired $t$-test, a bootstrap that resamples observations, or a test on fold-wise scores still relies on independence assumptions that the data violate. Whether an apparent improvement is ``statistically real'' then depends as much on these hidden assumptions as on the models.

Our starting point is the per-observation out-of-fold \emph{loss differential} $d_i = \ell(y_i,\hat f_A(x_i)) - \ell(y_i,\hat f_B(x_i))$ between two models $A$ and $B$. The mean of the $d_i$ estimates the expected gain of $A$ over $B$, so that the comparison becomes a problem of inference on the mean of a dependent random field. Two sources of dependence must be handled at once. (i)~Test-point dependence: the $d_i$ inherit the spatial and temporal autocorrelation of the residuals of both models. (ii)~Training dependence: all observations of a fold are predicted by the same fitted models, and the models of different folds are fitted on overlapping training sets. Existing procedures address at most one of these sources. Heteroskedasticity- and autocorrelation-consistent (HAC) variance estimators, as in the Diebold--Mariano test for forecasts~\cite{diebold1995}, and the parametric spatial prediction comparison test of Hering and Genton~\cite{hering2011} address (i) for fixed predictors, whereas the corrected resampled $t$-test of Nadeau and Bengio~\cite{nadeau2003} addresses (ii), but only for i.i.d. data.

\paragraph{Contributions.} The paper makes five contributions.
\begin{enumerate}
\item We formalise model comparison under spatiotemporal dependence as inference on the mean of a loss-differential field with correlated fold-level training effects. We show (Proposition~\ref{prop:pairedt}) that the paired $t$-test on fold scores underestimates the variance by exactly the average covariance between fold means, and we derive (Proposition~\ref{prop:bias}) the exact bias of a variance estimator combining test-sampling and training components, which identifies the conditions under which it is unbiased and the direction of its error otherwise.
\item We propose \emph{DA-P}, a dependence-aware test for two models evaluated by blocked cross-validation. It combines (a) a test-sampling variance computed from a separable correlation model fitted to the variograms of the fold-centred loss differentials; (b) a method-of-moments estimate of the training-induced fold variance, with an explicit assumption on the correlation of training effects; and (c) a Welch--Satterthwaite reference distribution. All steps use a single cross-validation run.
\item We attach to the test a three-way decision rule that declares a comparison \emph{indeterminate} when its significance depends on the unverifiable correlation of training effects, and we evaluate it, together with a range-based diagnostic, by simulation.
\item We assess conventional and dependence-aware procedures (paired $t$, i.i.d. bootstrap, corrected resampled $t$, block bootstrap, linear mixed model, Diebold--Mariano--HAC, a space--time version of the Hering--Genton test, a kernel variant DA-K, and DA-P) in simulation studies in which the truth is known exactly. These studies cover the sensitivity to the assumed training correlation, the robustness to anisotropy, covariance shape, heteroscedasticity, non-stationarity and irregular sampling, and gradient-boosting learners whose expected performances are equal by an exchangeability argument.
\item We apply all procedures to three public data sets (hourly weather in New York, hourly pedestrian counts in Melbourne and yield-monitor maize data in Argentina) under random, blocked and forward-chaining validation, with a correction for multiple comparisons and a ``seed-null'' diagnostic, and we derive practical recommendations. An open-source Python implementation reproduces all results.
\end{enumerate}

The emphasis of the paper is thus not on whether spatial or temporal validation is preferable, a question already addressed in the literature, but on the \emph{validity of the inference} drawn from a given validation design. In the terms of the call for this special issue, it contributes to measuring the trustworthiness of comparisons between machine learning models in real applications.

The rest of the paper is organised as follows. Section~\ref{sec:related} reviews related work. Section~\ref{sec:problem} sets up the problem and analyses conventional tests. Section~\ref{sec:method} presents DA-P and its decision rule. Sections~\ref{sec:sim1} and~\ref{sec:sim2} report the simulation studies, Section~\ref{sec:real} the case studies, and Section~\ref{sec:discussion} discusses recommendations and limitations.

\section{Related work}\label{sec:related}

\paragraph{Statistical comparison of learning algorithms.}
Dietterich~\cite{dietterich1998} showed that the paired $t$-test on resampled hold-out scores has an inflated type~I error rate, because training sets overlap, and proposed the $5\times2$ cross-validation test. Nadeau and Bengio~\cite{nadeau2003} derived the corrected resampled $t$-test, which inflates the variance of repeated train/test differences by the factor $1/J + n_{\text{test}}/n_{\text{train}}$. Bengio and Grandvalet~\cite{bengio2004} proved that no unbiased estimator of the variance of $K$-fold cross-validation exists, which is why every practical procedure relies on a working model or a heuristic. Bouckaert and Frank~\cite{bouckaert2004} studied the replicability of these tests, Dem\v{s}ar~\cite{demsar2006} and Benavoli et al.~\cite{benavoli2017} treated comparisons over multiple data sets, frequentist and Bayesian. More recently, Bayle et al.~\cite{bayle2020} established central limit theorems for $K$-fold cross-validation under algorithmic stability, and Bates et al.~\cite{bates2024} clarified that cross-validation estimates the average error of the learning \emph{procedure} rather than that of the fitted model, and proposed nested cross-validation intervals. Bouthillier et al.~\cite{bouthillier2021} quantified the sources of variance (data split, initialisation, hyperparameter search) in machine learning benchmarks, and Varoquaux~\cite{varoquaux2018} showed how wide cross-validation error bars are in small samples. All these works assume independent observations.

\paragraph{Validation under spatial and temporal dependence.}
In ecology, remote sensing and the geosciences, a large literature shows that random cross-validation is optimistic for spatially structured data and recommends spatial blocking, buffering or nearest-neighbour distance matching~\cite{roberts2017,valavi2019,lerest2014,meyer2018,ploton2020,schratz2019,linnenbrink2024}. Wadoux et al.~\cite{wadoux2021} argued that, when the goal is the accuracy of a map over a fixed region, design-based probability sampling is preferable to any form of spatial cross-validation. In time series, Bergmeir and Ben\'{i}tez~\cite{bergmeir2012} compared blocked and out-of-sample schemes and Bergmeir et al.~\cite{bergmeir2018} showed when standard cross-validation remains valid for autoregressive models. These works are concerned with the \emph{point estimate} of predictive performance and its bias; the uncertainty of the comparison between two models, which is the object of this paper, is rarely addressed.

\paragraph{Comparing predictions of dependent data.}
Inference on the difference in accuracy between two predictions of dependent data is not new. The Diebold--Mariano (DM) test~\cite{diebold1995}, with the small-sample correction of Harvey et al.~\cite{harvey1997}, applies a HAC (Newey--West~\cite{newey1987}) variance estimator to the loss-differential series of two forecasts. Hering and Genton~\cite{hering2011} proposed the spatial prediction comparison test, which models the covariance of a spatially correlated loss-differential field with a parametric variogram, and Gilleland et al.~\cite{gilleland2018} showed that the size of such tests is sensitive to serial dependence and to the correlation between competing forecasts. Watson et al.~\cite{watson2024} studied prediction and model evaluation for space--time data, with emphasis on how the validation design must match the prediction task. HAC estimation under cross-sectional dependence goes back to Conley~\cite{conley1999} and, for panels, to Driscoll and Kraay~\cite{driscoll1998}; the poor size of HAC tests in small effective samples motivated the fixed-$b$ theory of Kiefer and Vogelsang~\cite{kiefer2005}. Block and stationary bootstraps~\cite{kunsch1989,politis1994,lahiri2003} provide resampling alternatives.

\paragraph{Positioning of this work.}
In view of this literature, our contribution is not the idea of accounting for dependence in the loss differential, which the DM and Hering--Genton tests already do for \emph{fixed} predictors. It lies in four more specific elements. (i)~We compare \emph{learning procedures} that are refitted on each training fold of a blocked cross-validation, which adds a training-induced, fold-level variance component, correlated across folds, that is absent from forecast or map comparisons; we make its estimation and the conditions for its validity explicit (Propositions~\ref{prop:pairedt} and~\ref{prop:bias}). (ii)~DA-P combines this component with the test-sampling dependence in a single variance estimate computed from one cross-validation run. (iii)~The dependence of the test-sampling fluctuations is modelled on the variograms of the \emph{fold-centred} loss differentials, so that it is not confounded with the training effects, and its adequacy is examined under several misspecifications. (iv)~The procedure returns diagnostics and a three-way decision that declares a comparison indeterminate when the data contain too few effectively independent regions; the threshold of this rule is calibrated and validated by simulation. We include a space--time version of the Hering--Genton test among the competing procedures in all our experiments.

\paragraph{Trustworthy machine learning and this special issue.}
Recent contributions in this journal have advocated statistical tools for the robustness of machine learning models~\cite{giudici2024}, practical safeguards against overfitting~\cite{chicco2026}, and effect-size reporting beyond $p$-values in large data sets~\cite{caponero2025}. The present paper complements them by addressing the validity of the comparison step itself when observations are dependent.

\section{Problem formulation}\label{sec:problem}

\subsection{Setting and estimand}

We observe $n$ labelled examples $z_i=(x_i,y_i)$, $i=1,\dots,n$, each attached to a spatial location $s_i\in\mathbb{R}^p$ (typically $p=2$) and a time $t_i\in\mathbb{R}$; either component may be absent. Two learning procedures $A$ and $B$ (for instance, XGBoost and a random forest, or one learner with two feature sets) are to be compared using a loss $\ell$, such as the squared error or the logarithmic loss~\cite{gneiting2007}.

The observations are grouped into $G$ spatiotemporal blocks (for instance, square tiles crossed with time windows), and the blocks are allocated to $K$ folds $F_1,\dots,F_K$, with $n_k=|F_k|$. For each fold $k$, both learners are fitted on a training set $T_k\subseteq\{i\notin F_k\}$, from which observations within a buffer distance of the fold may be removed, and they are then used to predict the observations of $F_k$. With $k(i)$ the fold of observation $i$, the out-of-fold \emph{loss differential} is
\begin{equation}\label{eq:d}
d_i=\ell\bigl(y_i,\hat f_A^{(k(i))}(x_i)\bigr)-\ell\bigl(y_i,\hat f_B^{(k(i))}(x_i)\bigr),
\end{equation}
so that negative values favour $A$, and the cross-validated estimate of the gain is $\hat\Delta=\dbar=n^{-1}\sum_i d_i$.

\begin{definition}[Estimand]
The \emph{expected cross-validated loss differential} is $\Delta_n=\E[\hat\Delta]$, where the expectation is over the data-generating process and the randomness of the learners. It is the difference in expected loss between the two learning procedures, trained on about $n(1-1/K)$ observations, when predicting observations that are separated from the training data as prescribed by the blocking design.
\end{definition}

The estimand depends on the validation design, as stressed by Bates et al.~\cite{bates2024} for i.i.d. data and by the spatial validation literature~\cite{roberts2017,wadoux2021}: with random folds it describes interpolation, whereas with spatial blocks it describes extrapolation to new areas. We take the design as given, chosen to match the intended use of the model, and we focus on testing $H_0:\Delta_n=0$ against $H_1:\Delta_n\neq0$ and on confidence intervals for $\Delta_n$.

\subsection{A working model for loss differentials}

The difficulty is that the $d_i$ are dependent for two distinct reasons, which we make explicit with the following decomposition.

\begin{definition}[Working model]\label{def:model}
\begin{equation}\label{eq:model}
d_i=\Delta_n+\eta_{k(i)}+u_i ,
\end{equation}
where $(\eta_1,\dots,\eta_K)$ are exchangeable fold effects with $\Var(\eta_k)=\tau^2$ and $\operatorname{Corr}(\eta_k,\eta_l)=\rho$ for $k\neq l$, independent of $u$, and $u=(u_i)$ is a zero-mean, second-order stationary field with $\Cov(u_i,u_j)=C(s_i-s_j,t_i-t_j)$.
\end{definition}

The fold effect $\eta_k$ captures the variability induced by fitting the two learners on the training set $T_k$: an unlucky training sample, a different random seed or a different tree structure shifts the losses of all observations of the fold in the same direction. Because the training sets of different folds share a fraction $(K-2)/(K-1)$ of their observations, the fold effects are positively correlated. The field $u$ captures the test-sampling variability; it inherits the spatial and temporal autocorrelation of the residuals of both models, which is substantial whenever both models miss a spatially structured signal. Under~\eqref{eq:model},
\begin{equation}\label{eq:var}
\Var(\hat\Delta)=\underbrace{\frac{1}{n^2}\sum_{i=1}^n\sum_{j=1}^n C(s_i-s_j,t_i-t_j)}_{V_u}+\underbrace{\tau^2\Bigl[\sum_k\frac{n_k^2}{n^2}+\rho\sum_{k\neq l}\frac{n_kn_l}{n^2}\Bigr]}_{V_\eta}.
\end{equation}
With equal fold sizes, $V_\eta=\tau^2[1/K+\rho(1-1/K)]$. Writing $V_u=\sigma_u^2/n_{\mathrm{eff}}$ defines the \emph{effective sample size} $n_{\mathrm{eff}}=n/[1+(n-1)\bar\rho_u]$, where $\bar\rho_u$ is the average correlation between distinct $u_i$. Under moderate dependence $n_{\mathrm{eff}}$ can be orders of magnitude smaller than $n$.

\subsection{Why conventional procedures fail}

\paragraph{Procedures that assume independent observations.}
The i.i.d. bootstrap of the $d_i$, the $z$-test on the $d_i$ and McNemar-type tests estimate $\sigma^2/n$ and therefore ignore all off-diagonal terms of~\eqref{eq:var}. Their variance is too small by the design-effect factor $n/n_{\mathrm{eff}}$, and it misses $V_\eta$ even in the absence of spatial dependence.

\paragraph{The paired $t$-test on fold scores.}
The most common procedure computes the $K$ fold means $m_k=n_k^{-1}\sum_{i\in F_k}d_i$ and refers $\bar m/(s_m/\sqrt K)$ to a $t_{K-1}$ distribution, with $s_m^2$ the sample variance of the $m_k$. Its weakness is captured by the following identity.

\begin{proposition}\label{prop:pairedt}
Let $m_1,\dots,m_K$ have a common mean, average variance $\bar v=K^{-1}\sum_k\Var(m_k)$ and average pairwise covariance $\bar c=[K(K-1)]^{-1}\sum_{k\neq l}\Cov(m_k,m_l)$. Then
\[
\E\Bigl[\frac{s_m^2}{K}\Bigr]-\Var(\bar m)=-\bar c .
\]
\end{proposition}

\begin{proof}
$\Var(\bar m)=K^{-2}\bigl[\sum_k\Var(m_k)+\sum_{k\neq l}\Cov(m_k,m_l)\bigr]=[\bar v+(K-1)\bar c]/K$. Since the means are equal, $\E\sum_k(m_k-\bar m)^2=\sum_k\Var(m_k)-K\Var(\bar m)=(K-1)(\bar v-\bar c)$, hence $\E[s_m^2]/K=(\bar v-\bar c)/K$. Subtracting gives the result.
\end{proof}

The paired $t$-test is therefore valid only if the fold means are uncorrelated. Under random $K$-fold cross-validation of dependent data, every fold is spread over the whole domain, the folds share the same spatial and temporal fluctuations, and $\bar c$ is close to $\Var(\bar m)$ itself, so that the test estimates only the small idiosyncratic part of the variance. Overlapping training sets add $\rho\tau^2>0$ to $\bar c$, which is the mechanism identified by Dietterich~\cite{dietterich1998} and corrected by Nadeau and Bengio~\cite{nadeau2003} for i.i.d. data. Blocking reduces the spatial part of $\bar c$, but adjacent blocks of different folds remain correlated, and with $K=5$ the test has only four degrees of freedom.

\paragraph{Cluster methods.}
The block bootstrap and the random-intercept linear mixed model (LMM) treat the $G$ blocks as independent clusters. They capture all within-block covariance, including the fold effects of~\eqref{eq:model}, but they ignore the covariance between neighbouring blocks, which is non-negligible when the dependence range is comparable with the block size, as well as the correlation $\rho$ of fold effects.

\paragraph{Parametric spatial comparison.}
The test of Hering and Genton~\cite{hering2011}, which we extend to space--time with a separable exponential covariance fitted to the spatial and temporal variograms of the centred differentials, estimates $V_u$ from a parametric model. It is designed for fixed predictors and therefore omits $V_\eta$, and its validity depends on the adequacy of the covariance model.

\paragraph{Time-only HAC.}
Applying the Diebold--Mariano test with a Newey--West variance separately to the series of each site, as is often done for panels of forecasts, accounts for temporal dependence but ignores both the correlation between sites at the same time and the fold effects.

\section{Dependence-aware HAC inference}\label{sec:method}

DA-HAC estimates the two components of~\eqref{eq:var} from the loss differentials of a single blocked cross-validation run. Figure~\ref{fig:architecture} gives an overview of the complete procedure: blocked cross-validation produces the loss differentials (stage~1); two parallel branches estimate the test-sampling variance, through variogram-calibrated kernel weights, and the training-induced variance, from the dispersion of fold means (stage~2); the two are combined into a variance, a reference distribution and diagnostics (stage~3). The steps are summarised in Algorithm~\ref{alg:dahac} and detailed below.

\begin{figure}[t]
\centering
\resizebox{\textwidth}{!}{% Architecture of the proposed comparison procedure (included by main.tex).
% Reading order: stage 1 left to right, stage 2 right to left, stage 3 left to right.
\begin{tikzpicture}[
    font=\footnotesize,
    >=Stealth,
    node distance=5mm and 5mm,
    box/.style={draw=#1!75!black, fill=#1!10, rounded corners=2pt, align=center,
                minimum height=13mm, text width=29mm, inner sep=2pt},
    box/.default=gray,
    lane/.style={draw=#1!70!black, dashed, rounded corners=4pt, inner sep=2.5mm},
    arr/.style={->, thick, draw=black!65},
    lbl/.style={font=\scriptsize\bfseries, text=#1!70!black, anchor=south west, inner sep=1pt},
]
\definecolor{cpipe}{HTML}{52514E}
\definecolor{ctest}{HTML}{2A78D6}
\definecolor{ctrain}{HTML}{EB6834}
\definecolor{cout}{HTML}{1BAF7A}

% ---------------- Stage 1: blocked cross-validation ----------------
\node[box=cpipe] (data) {\textbf{Data}\\[1pt] $(x_i,y_i)$ observed at\\ location $s_i$, time $t_i$};
\node[box=cpipe, right=of data] (block) {\textbf{Blocking}\\[1pt] $G$ space--time blocks\\ $\to K$ folds, buffer};
\node[box=cpipe, right=of block] (fit) {\textbf{Refit $A$ and $B$}\\[1pt] on each training set\\ $T_k$, predict fold $F_k$};
\node[box=cpipe, right=of fit] (diff) {\textbf{Loss differentials}\\[1pt] $d_i=\ell_A(i)-\ell_B(i)$\\ estimate $\hat\Delta=\bar d$};
\draw[arr] (data) -- (block);
\draw[arr] (block) -- (fit);
\draw[arr] (fit) -- (diff);
\begin{scope}[on background layer]
  \node[lane=cpipe, fit=(data)(diff)] (L1) {};
\end{scope}
\node[lbl=cpipe] at (L1.north west) {1\quad Blocked cross-validation (defines the estimand $\Delta_n$)};

% ---------------- Stage 2: DA-P variance (right to left) ----------------
\node[box=ctest, below=11mm of diff] (center) {\textbf{Fold-centring}\\[1pt] fold means $m_k$\\ $\tilde u_i=d_i-m_{k(i)}$};
\node[box=ctest, left=of center] (vario) {\textbf{Variograms of $\tilde u$}\\[1pt] in space and in time\\ $\to$ ranges $\hat r_S,\hat r_T$};
\node[box=ctest, left=of vario] (bw) {\textbf{Correlation model}\\[1pt] exponential + nugget\\ $\hat g_{ij}=\hat g_S(h_{ij})\,\hat g_T(\ell_{ij})$};
\node[box=ctest, left=of bw] (vu) {\textbf{Test-sampling var.}\\[1pt] $\hat V_u=\frac{\hat\sigma^2}{n^2}\sum_{i,j}\hat g_{ij}$\\ and per fold $\hat V_k$};
\draw[arr] (diff) -- (center);
\draw[arr] (center) -- (vario);
\draw[arr] (vario) -- (bw);
\draw[arr] (bw) -- (vu);

\node[box=ctrain, below=of center] (fm) {\textbf{Between-fold}\\[1pt] variance of the\\ fold means: $s_m^2$};
\node[box=ctrain, left=of fm, text width=63mm] (tau) {\textbf{Training variance} (method of moments)\\[1pt] $\hat\tau^2=\max\bigl\{0,\ s_m^2-K^{-1}\textstyle\sum_k\hat V_k\bigr\}$};
\node[box=ctrain, left=of tau] (vt) {\textbf{Nadeau--Bengio}\\[1pt] $\zeta_K\hat\tau^2$\\ $\zeta_K=\tfrac1K+\tfrac1{K-1}$};
\draw[arr] (center) -- (fm);
\draw[arr] (fm) -- (tau);
\draw[arr] (tau) -- (vt);
\draw[arr, dashed] ([xshift=-3mm]vu.south east) -- node[midway, below left=-1pt, font=\scriptsize]{$\hat V_k$} ([xshift=6mm]tau.north west);

\begin{scope}[on background layer]
  \node[lane=ctest, fit=(center)(vu)(fm)(vt)] (L2) {};
\end{scope}
\node[lbl=ctest] at (L2.north west) {2\quad DA-P variance estimation};

% ---------------- Stage 3: inference and diagnostics ----------------
\node[box=cout, below=11mm of vt.south west, anchor=north west, text width=40mm] (v) {\textbf{Total variance}\\[1pt] $\hat V=\hat V_u+\zeta_K\hat\tau^2$};
\node[box=cout, right=of v, text width=40mm] (df) {\textbf{Satterthwaite d.f.\ $\hat\nu$}\\[1pt] from $\hat n_{\mathrm{eff}}$, $\textstyle\prod_q 3L_q/(2\hat r_q)$\\ and $K-1$; range ratio $\hat r/L$};
\node[box=cout, right=of df, text width=40mm] (out) {\textbf{Output}\\[1pt] $T=\hat\Delta/\sqrt{\hat V}\sim t_{\hat\nu}$: $p$-value, CI\\ decision: reject / not / indet.};
\draw[arr] (vt.south) -- (vt.south |- v.north);
\draw[arr] (vu.west) -- ++(-5mm,0) |- (v.west);
\draw[arr] (v) -- (df);
\draw[arr] (df) -- (out);

\begin{scope}[on background layer]
  \node[lane=cout, fit=(v)(out)] (L3) {};
\end{scope}
\node[lbl=cout, anchor=south east] at (L3.north east) {3\quad Inference and diagnostics};
\end{tikzpicture}
}
\caption{Architecture of the proposed comparison procedure. Stage~1 (grey) is a standard blocked cross-validation that yields the per-observation loss differentials $d_i$ and defines the estimand. Stage~2 (DA-HAC) estimates the test-sampling variance $\hat V_u$ (blue branch, read from right to left: fold-centring, variograms, bandwidths, product-Bartlett kernel) and the training-induced variance $\zeta_K\hat\tau^2$ (orange branch). Stage~3 (green) combines them into the total variance, the Satterthwaite degrees of freedom and the test, together with the diagnostics used in Section~\ref{sec:discussion}.}\label{fig:architecture}
\end{figure}

\subsection{Fold-centred differentials}

Let $m_k$ be the mean of the $d_i$ in fold $k$ and $\tilde u_i=d_i-m_{k(i)}$. Under~\eqref{eq:model}, $\tilde u_i=u_i-\bar u_{k(i)}$ is free of the fold effects, so that the dependence structure of $u$ can be studied without confusion with the training-induced component.

\subsection{Variogram-calibrated bandwidths}

The key tuning parameter of any HAC estimator is the bandwidth, i.e.\ the distance beyond which covariances are neglected. Rules of thumb developed for time series~\cite{newey1987,andrews1991} do not transfer to space, and the dependence range of the loss differential is in general different from that of the response: it is the range of the \emph{residual} structure that both models fail to capture. We therefore estimate it from $\tilde u$ directly. The directional empirical semivariograms
\begin{align*}
\hat\gamma_S(h)&=\frac{1}{2|N_S(h)|}\sum_{(i,j)\in N_S(h)}(\tilde u_i-\tilde u_j)^2,& N_S(h)&=\{(i,j):t_i=t_j,\ \|s_i-s_j\|\in I_h\},\\
\hat\gamma_T(\ell)&=\frac{1}{2|N_T(\ell)|}\sum_{(i,j)\in N_T(\ell)}(\tilde u_i-\tilde u_j)^2,& N_T(\ell)&=\{(i,j):s_i=s_j,\ |t_i-t_j|\in I_\ell\},
\end{align*}
are computed over distance and lag bins $I_h$, $I_\ell$ up to a maximum lag (by default half the spatial extent and a quarter of the temporal extent). For purely spatial data, $N_S(h)$ uses all pairs. To each we fit the exponential model $\gamma(h)=c_0+c_1\{1-\exp(-h/a)\}$ by the weighted least squares of Cressie~\cite{cressie1985}, with weights $|N(h)|/\gamma(h)^2$. The \emph{practical range} $r=3a$ is the distance at which the correlation of the structured component falls to $5\%$; it is set to zero when the structured part explains less than $5\%$ of the sill ($c_1<0.05(c_0+c_1)$). The spatial and temporal bandwidths are
\[
b_S=\kappa\,\hat r_S,\qquad b_T=\kappa\,\hat r_T ,
\]
with $\kappa=1$ chosen in a pilot study on separate random seeds (Section~\ref{sec:sim1}). We also report the \emph{range ratio} $\max_q b_q/L_q$, the largest bandwidth relative to the extent of the domain along the same axis, as a diagnostic: when it approaches one, the sample contains only a handful of effectively independent regions, and no procedure can estimate the variance of $\hat\Delta$ reliably from it.

\subsection{Kernel estimate of the test-sampling variance}

We use the product-Bartlett kernel
\begin{equation}\label{eq:kernel}
w_{ij}=\prod_{q=1}^{p}\Bigl(1-\frac{|s_{iq}-s_{jq}|}{b_S}\Bigr)_+\Bigl(1-\frac{|t_i-t_j|}{b_T}\Bigr)_+ ,
\end{equation}
where a zero bandwidth keeps only pairs with identical coordinates on that axis. Each factor is a one-dimensional Bartlett (triangular) kernel, whose Fourier transform is non-negative; the product is therefore a positive semi-definite kernel on $\mathbb{R}^{p+1}$, which guarantees non-negative variance estimates. (A Bartlett kernel applied to the Euclidean distance in two or more dimensions does not have this property.) The test-sampling variance and its within-fold analogues are
\begin{equation}\label{eq:vu}
\hat V_u=\frac{1}{n^2}\sum_{i,j}w_{ij}\tilde u_i\tilde u_j,\qquad
\hat V_k=\frac{1}{n_k^2}\sum_{i,j\in F_k}w_{ij}\tilde u_i\tilde u_j .
\end{equation}
Only pairs with $w_{ij}>0$ contribute; they are enumerated with a $k$-d tree in the Chebyshev metric after scaling each coordinate by its bandwidth, at a cost proportional to the number of such pairs.

\subsection{Training-induced variance}

Under~\eqref{eq:model}, the between-fold variance of fold means exceeds their within-fold test-sampling variance by the variance of the fold effects. We estimate it by the method of moments,
\begin{equation}\label{eq:tau}
\hat\tau^2=\max\Bigl\{0,\ s_m^2-\frac1K\sum_{k=1}^K\hat V_k\Bigr\}.
\end{equation}
Since $\E[s_m^2]$ subtracts the covariance between fold effects (Proposition~\ref{prop:pairedt}), $\hat\tau^2$ targets $\tau^2(1-\rho)$ rather than $\tau^2$. The correlation $\rho$ is not identifiable from one cross-validation run~\cite{bengio2004}. For an assumed value $\rho_0$, the training component of~\eqref{eq:var} is $V_\eta=\tau^2[1/K+\rho_0(1-1/K)]=\zeta_K(\rho_0)\,\tau^2(1-\rho_0)$, with
\begin{equation}\label{eq:zeta}
\zeta_K(\rho_0)=\frac{1/K+\rho_0(1-1/K)}{1-\rho_0},\qquad \zeta_K(1/K)=\frac1K+\frac1{K-1}.
\end{equation}
Following Nadeau and Bengio~\cite{nadeau2003}, our default is the test fraction $\rho_0=n_{\text{test}}/n=1/K$, for which $\zeta_K$ is exactly their inflation $1/J+n_{\text{test}}/n_{\text{train}}$ with $J=K$ splits. The final variance estimate is
\begin{equation}\label{eq:vtotal}
\hat V=\hat V_u+\zeta_K(\rho_0)\,\hat\tau^2 .
\end{equation}

\subsection{What the variance estimator targets}

The following proposition makes explicit the conditions under which~\eqref{eq:vtotal} is unbiased, and the direction of its bias when they fail. We consider the untruncated estimator $\hat\tau^2_{\mathrm{raw}}=s_m^2-K^{-1}\sum_k\hat V_k$ and $\hat V_{\mathrm{raw}}=\hat V_u+\zeta_K(\rho_0)\hat\tau^2_{\mathrm{raw}}$.

\begin{proposition}\label{prop:bias}
Under the working model~\eqref{eq:model} with equal fold sizes, let $\bar m^u_k$ be the fold means of $u$, $v_u=K^{-1}\sum_k\Var(\bar m^u_k)$, $c_u=[K(K-1)]^{-1}\sum_{k\neq l}\Cov(\bar m^u_k,\bar m^u_l)$, and let $\beta_u=\E[\hat V_u]-V_u$ and $\bar\beta=K^{-1}\sum_k\{\E[\hat V_k]-\Var(\bar m^u_k)\}$ be the biases of the kernel estimators. Then
\begin{equation}\label{eq:bias}
\E[\hat V_{\mathrm{raw}}]-\Var(\hat\Delta)
=\underbrace{\beta_u-\zeta_K(\rho_0)\,(c_u+\bar\beta)}_{\text{test-sampling part}}
+\underbrace{\tau^2\Bigl[(1-\rho)\,\zeta_K(\rho_0)-\tfrac1K-\rho\bigl(1-\tfrac1K\bigr)\Bigr]}_{\text{training part}} .
\end{equation}
In particular, $\hat V_{\mathrm{raw}}$ is unbiased if (i) $\rho_0=\rho$, (ii) fold means of $u$ are uncorrelated across folds ($c_u=0$) and (iii) the kernel estimators are unbiased ($\beta_u=\bar\beta=0$). The training part is negative, i.e.\ the variance is underestimated, if and only if $\rho>\rho_0$.
\end{proposition}

\begin{proof}
The fold means are $m_k=\Delta_n+\eta_k+\bar m^u_k$ with $\eta$ independent of $u$, so their average variance is $\tau^2+v_u$ and their average covariance $\rho\tau^2+c_u$. By the identity in the proof of Proposition~\ref{prop:pairedt}, $\E[s_m^2]=(1-\rho)\tau^2+v_u-c_u$, and by definition $\E[K^{-1}\sum_k\hat V_k]=v_u+\bar\beta$, hence $\E[\hat\tau^2_{\mathrm{raw}}]=(1-\rho)\tau^2-c_u-\bar\beta$. With equal fold sizes, \eqref{eq:var} reads $\Var(\hat\Delta)=V_u+\tau^2[1/K+\rho(1-1/K)]$ and $\E[\hat V_u]=V_u+\beta_u$. Substituting gives~\eqref{eq:bias}. Writing $g(r)=1/K+r(1-1/K)$, the training part is $\tau^2[(1-\rho)g(\rho_0)/(1-\rho_0)-g(\rho)]$, which vanishes at $\rho=\rho_0$ and is strictly decreasing in $\rho$, with derivative $-\tau^2[g(\rho_0)/(1-\rho_0)+1-1/K]<0$.
\end{proof}

Proposition~\ref{prop:bias} separates three sources of error, each of which is evaluated empirically in Section~\ref{sec:sim1}. \emph{Training correlation.} For $K=5$ and $\rho_0=1/K$, the training part equals $\tau^2(0.25-1.25\rho)$: assuming $\rho_0=0.2$ when $\rho=0.5$ leaves $0.375\tau^2$ of the $0.6\tau^2$ training variance unaccounted for, whereas $\rho<\rho_0$ errs on the conservative side. The sensitivity of the test to this choice is reported in Section~\ref{sec:sim1rho}. \emph{Cross-fold covariance.} $c_u$ is the covariance between test-sampling fluctuations of different folds; it is reduced by allocating whole blocks to folds and by buffers, and vanishes when blocks are larger than the dependence range. Because it enters with a negative sign, residual cross-fold dependence makes the test anti-conservative. \emph{Kernel bias.} Fold-centring removes the fold-level average of $u$ from $\tilde u$, and the Bartlett kernel down-weights covariances near the bandwidth; both make $\beta_u$ and $\bar\beta$ negative when the dependence range is not small relative to the blocks and to the domain. The term $-\zeta_K\bar\beta$ partially offsets $\beta_u$: the part of the fold-level variance missed by the within-fold estimators is recovered by $\hat\tau^2$. Finally, truncation at zero, $\hat\tau^2=\max\{0,\hat\tau^2_{\mathrm{raw}}\}$, can only increase the variance, so $\E[\hat V]\ge\E[\hat V_{\mathrm{raw}}]$; its effect on the reference distribution is discussed next.

\subsection{Reference distribution}

HAC tests are known to over-reject when the effective sample size is small~\cite{kiefer2005}, and $\hat\tau^2$ rests on only $K-1$ degrees of freedom. We approximate the distribution of $\hat V$ by a scaled $\chi^2$ with Welch--Satterthwaite~\cite{satterthwaite1946} degrees of freedom
\begin{equation}\label{eq:df}
\hat\nu=\frac{\hat V^2}{\hat V_u^2/\nu_u+(\zeta_K\hat\tau^2)^2/(K-1)},\qquad
\nu_u=\min\Bigl\{\hat n_{\mathrm{eff}}-1,\ \prod_{q}\frac{3L_q}{2b_q}\Bigr\},
\end{equation}
with $\hat n_{\mathrm{eff}}=n^{-1}\sum_i\tilde u_i^2/\hat V_u$ truncated to $[2,n]$, where the product runs over the spatial and temporal axes with a positive bandwidth, and $L_q$ is the extent of the domain along axis $q$. The second term of $\nu_u$ is the classical approximation to the degrees of freedom of a kernel spectral estimate~\cite{andrews1991}: for a lag window $k$ of bandwidth $b$ on a record of length $L$, $\Var(\hat V)/V^2\approx2(b/L)\int k^2$, with $\int k^2=2/3$ for the Bartlett window, and the factors multiply for a product kernel. It caps the degrees of freedom when the bandwidth is not small compared with the domain, the regime in which HAC tests are most anti-conservative. The test statistic $T=\hat\Delta/\sqrt{\hat V}$ is referred to a Student $t_{\hat\nu}$ distribution, and the $(1-\alpha)$ confidence interval is $\hat\Delta\pm t_{\hat\nu,1-\alpha/2}\sqrt{\hat V}$.

When $\hat\tau^2_{\mathrm{raw}}\le0$ the training term vanishes and $\hat\nu=\nu_u$, which may be large even though only $K$ folds inform the training component. Truncation occurs when the observed dispersion of fold means is no larger than its predicted test-sampling part, that is, when the data show no evidence of a training effect; the variance itself is then not underestimated on average (Proposition~\ref{prop:bias}), but its sampling variability is understated. Section~\ref{sec:sim1} reports the rejection rate conditionally on truncation to check that this boundary case does not cause over-rejection.

\subsection{A three-way decision rule}\label{sec:rule}

Proposition~\ref{prop:bias} shows that no correction of the variance can be exact when the dependence range is comparable to the blocks and to the study domain: $c_u$, $\beta_u$ and $\bar\beta$ are then all non-negligible and are not identifiable from one sample. We therefore attach to DA-HAC an explicit decision rule with three outcomes:
\begin{equation}\label{eq:rule}
\text{decision}=\begin{cases}
\text{indeterminate} & \text{if } \max_q b_q/L_q>r^\ast,\\
\text{reject } H_0 & \text{if } \max_q b_q/L_q\le r^\ast\text{ and } p<\alpha,\\
\text{do not reject } H_0 & \text{otherwise,}
\end{cases}
\end{equation}
where $b_q/L_q$ is the estimated bandwidth relative to the extent of the domain along axis $q$ (the range ratio). An ``indeterminate'' outcome states that the sample contains too few effectively independent regions for the comparison to be resolved, whatever the observed difference. The threshold $r^\ast=\rstar$ was chosen in a pilot study on separate random seeds as the largest value for which the conditional rejection rate of a true null stayed within the Monte Carlo band of the nominal level, and was then validated on new seeds and on the other experiments (Section~\ref{sec:sim1rule}).

\begin{algorithm}[t]
\caption{DA-HAC test for $H_0:\Delta_n=0$}\label{alg:dahac}
\begin{algorithmic}[1]
\Require loss differentials $d_i$, folds $k(i)$, coordinates $(s_i,t_i)$, level $\alpha$, multiplier $\kappa$ (default 1)
\State $m_k\gets$ mean of $d_i$ over $F_k$; \ $\tilde u_i\gets d_i-m_{k(i)}$; \ $\hat\Delta\gets n^{-1}\sum_i d_i$
\State Fit exponential models to the spatial and temporal empirical variograms of $\tilde u$; \ $b_S\gets\kappa\hat r_S$, $b_T\gets\kappa\hat r_T$
\State Enumerate pairs with $w_{ij}>0$ under the product-Bartlett kernel~\eqref{eq:kernel}
\State Compute $\hat V_u$ and $\hat V_1,\dots,\hat V_K$ by~\eqref{eq:vu}
\State $\hat\tau^2\gets\max\{0,s_m^2-K^{-1}\sum_k\hat V_k\}$; \ $\hat V\gets\hat V_u+\zeta_K(\rho_0)\hat\tau^2$ \Comment{default $\rho_0=1/K$}
\State Compute $\hat n_{\mathrm{eff}}$, $\nu_u$ and $\hat\nu$ by~\eqref{eq:df}, and the range ratio $\max_q b_q/L_q$
\State $p\gets2\,\Pr(t_{\hat\nu}>|\hat\Delta|/\sqrt{\hat V})$; interval $\hat\Delta\pm t_{\hat\nu,1-\alpha/2}\sqrt{\hat V}$
\State \Return $p$, the interval, the diagnostics and the decision~\eqref{eq:rule}: indeterminate if $\max_q b_q/L_q>r^\ast$
\end{algorithmic}
\end{algorithm}

\subsection{Properties}

\begin{remark}[Special cases]
If $b_S=b_T=0$ and $\hat\tau^2=0$, then $\hat V=n^{-2}\sum_i\tilde u_i^2$ and DA-HAC reduces to the one-sample $t$-test on the $d_i$. If the data have no spatial (or no temporal) coordinate, the corresponding factor of~\eqref{eq:kernel} is omitted. Because the kernel is positive semi-definite, $\hat V\ge0$ always.
\end{remark}

\begin{remark}[Consistency]
For fixed bandwidths, $\hat V_u$ is a Conley-type spatial HAC estimator~\cite{conley1999}. Under increasing-domain asymptotics with mixing conditions and bandwidths that grow with, but more slowly than, the size of the domain, such estimators are consistent~\cite{conley1999,andrews1991}. Our bandwidths are data-dependent and our variance includes a component, $\hat\tau^2$, estimated from only $K$ folds; we therefore do not claim exact asymptotic validity, and instead evaluate the finite-sample behaviour of the whole procedure by simulation with known truth.
\end{remark}

\begin{remark}[Scope]
DA-HAC requires only the per-observation losses of the two models and their coordinates; it is agnostic to the learners, to the loss and to the blocking scheme. The same construction applies to any additive performance measure (squared error, absolute error, log loss, Brier score, 0--1 loss). Non-additive measures such as the AUC or the F1-score would require a linearisation (influence function) first.
\end{remark}

\section{Simulation study I: controlled loss-differential fields}\label{sec:sim1}

The first study isolates the inferential problem from the learning problem. Prediction errors are simulated directly, so that the true expected loss differential is known exactly and the dependence of the $d_i$ is controlled.

\subsection{Design}

\paragraph{Fields.} Observations lie on a $12\times12$ grid of sites observed at $24$ time points ($n=3456$). All fields are separable Gaussian fields with unit variance, exponential spatial correlation $\exp(-h/a)$ and AR(1) temporal dependence with coefficient $\phi$. Four dependence levels are considered: \emph{none} ($a=0$, $\phi=0$), \emph{weak} ($a=1$, $\phi=0.3$), \emph{moderate} ($a=2$, $\phi=0.6$) and \emph{strong} ($a=4$, $\phi=0.8$). The corresponding practical ranges are $3a=0,3,6,12$ grid units in space and $-3/\log\phi=0,2.5,5.9,13.4$ time steps; at the strong level the range is as large as the domain.

\paragraph{Losses.} With $\varepsilon$, $e_A$ and $e_B$ three independent fields, the response noise is $\varepsilon$, model $A$ has errors $e_A$ and model $B$ has errors $\sqrt{1+\delta}\,e_B$. Under squared loss,
\[
d_i=(\varepsilon_i-e_{A,i})^2-(\varepsilon_i-\sqrt{1+\delta}\,e_{B,i})^2+\eta_{k(i)},\qquad \E[d_i]=-\delta ,
\]
so that $\Delta_n=-\delta$ exactly. The fold effects $\eta_k=\sigma_\eta(\sqrt{\rho}\,g+\sqrt{1-\rho}\,h_k)$, with $g,h_1,\dots,h_K$ i.i.d. standard normal, $\sigma_\eta=0.05$ and $\rho=0.25$, mimic the training-induced variability of~\eqref{eq:model}, including its correlation across folds. We use $\delta=0$ (null) and $\delta\in\{0.1,0.2,0.4\}$ (alternatives).

\paragraph{Validation and procedures.} The domain is cut into $4\times4$-site tiles crossed with $6$-step time windows ($G=36$ blocks), which are allocated to $K=5$ folds with balanced sizes. We compare the paired $t$-test on fold means, the i.i.d. bootstrap, the block bootstrap and a random-intercept LMM (both with the $36$ blocks as clusters), the Diebold--Mariano test with a site-wise Newey--West variance (DM--HAC), and DA-HAC. Three ablations of DA-HAC are included: without the training term ($\hat\tau^2=0$), with a normal reference distribution, and with a fixed bandwidth equal to the block size instead of the variogram-calibrated one. Bootstraps use $999$ resamples. There are $1000$ replicates per dependence level under the null and $300$ under each alternative, so that the Monte Carlo standard error of a rejection rate of $5\%$ is $0.7$ percentage points.

\paragraph{Pilot.} The bandwidth multiplier $\kappa$ and the reference distribution were chosen in a pilot study with $300$ null replicates per level, run on random seeds disjoint from those of the main study. Rejection rates of DA-HAC varied by at most $2.7$ percentage points across $\kappa\in\{0.5,1,1.5,2\}$ at any level, with no systematic trend, and we retained $\kappa=1$. Capping the degrees of freedom by the kernel term of~\eqref{eq:df} reduced the rejection rate at the strong level from $17.7\%$ to $15.3\%$ without affecting the other levels by more than $1.4$ points, and was retained. The pilot results are included in the code repository.

%%SIM1RESULTS%%

\section{Simulation study II: learners on spatio-temporal data}\label{sec:sim2}

The second study reproduces the full pipeline of an applied comparison: data generation, model fitting by cross-validation and testing. The difficulty in such a study is to know the truth. We obtain an exact null by exchangeability.

\subsection{Design}

\paragraph{Data.} On a $10\times10$ grid observed at $20$ time points ($n=2000$), three covariates $x_1,x_2,x_3$ and a latent residual $w$ are independent spatio-temporal fields with the dependence levels of Section~\ref{sec:sim1}. The response is
\[
y=g(x_1)+(1-\gamma)\,g(x_2)+x_3+w+\epsilon,\qquad g(v)=\sin(2v)+v,\qquad \epsilon\sim N(0,0.5^2)\ \text{i.i.d.}
\]
\paragraph{Models.} Model $A$ is XGBoost~\cite{chen2016} ($100$ trees of depth $4$, learning rate $0.1$, row subsampling $0.8$) fitted on $(x_1,x_3,s,t)$; model $B$ is the same learner fitted on $(x_2,x_3,s,t)$. Including the coordinates lets both models exploit the spatial and temporal structure of $w$, as is common in practice.

\paragraph{Exact null by exchangeability.} When $\gamma=0$, the joint distribution of the data is invariant under the exchange of $x_1$ and $x_2$, which exchanges the roles of $A$ and $B$. Hence $\E[d_i]=0$ for every observation and every validation design: the two models have exactly the same expected performance, although in any given sample one of them will appear better. When $\gamma>0$, $x_1$ carries more signal than $x_2$ and $A$ is truly better; we use $\gamma\in\{0.1,0.2,0.3\}$.

\paragraph{Validation and procedures.} Two pipelines are applied to each data set. The \emph{conventional} pipeline uses random $5$-fold cross-validation with the paired $t$-test on fold means and the i.i.d. bootstrap, and $5\times5$ repeated random cross-validation with the corrected resampled $t$-test~\cite{nadeau2003}. The \emph{dependence-aware} pipeline uses blocked $5$-fold cross-validation ($5\times5$-site tiles crossed with $4$-step windows, $G=20$ blocks, and a buffer that removes from the training set observations within one site and one time step of the test fold) with the procedures of Section~\ref{sec:sim1}. There are $300$ replicates per dependence level under the null and $150$ under each alternative.

%%SIM2RESULTS%%

\section{Case studies}\label{sec:real}

\subsection{Data and protocol}

We use three public data sets obtained from the Rdatasets collection~\cite{rdatasets}.

\paragraph{NYC weather.} Hourly weather records at the three New York City airports (EWR, JFK, LGA) in 2013, distributed with the \texttt{nycflights13} package~\cite{nycflights13} ($n=26{,}110$ after removing five records with missing temperature, humidity or wind speed). The task is to classify whether precipitation was recorded in the hour, from temperature, dew point, humidity, wind direction and speed, pressure, month, hour and airport; $6.7\%$ of hours are positive. The loss is the logarithmic loss.

\paragraph{Melbourne pedestrians.} Hourly pedestrian counts at four sensors in Melbourne in 2015--2016, distributed with the \texttt{tsibble} package~\cite{wang2020tsibble} ($n=66{,}037$). The task is to predict $\log(1+\text{count})$ from hour, day of week, month, day of year, year and sensor, under squared loss.

\paragraph{Las Rosas maize.} Yield-monitor data from a maize field in Las Rosas, Argentina, in 1999 and 2001~\cite{anselin2004}, distributed with the \texttt{agridat} package ($n=3443$, field of about $660\times380$\,m). The task is to predict yield from the nitrogen rate, brightness value, topographic zone, year and coordinates, under squared loss.

For the two hourly data sets, blocks are two-week windows ($26$ and $53$ blocks); as the airports lie within about $35$\,km of one another and the sensors within about $2$\,km, they are treated as co-located, so that all sites observed at the same time may be correlated, and a $24$-hour buffer is removed from the training set on each side of a test block. For the maize data, blocks are $75\times75$\,m tiles ($35$ non-empty blocks) with a $10$\,m buffer. We compare XGBoost ($300$ trees of depth $6$, learning rate $0.05$, row and column subsampling $0.8$) with a random forest~\cite{breiman2001} ($200$ trees, minimum leaf size $5$, $50\%$ of features per split) and with a linear baseline (logistic or ridge regression on standardised, mean-imputed features). All models use scikit-learn~\cite{pedregosa2011} or the XGBoost library with default settings otherwise; no hyperparameter was tuned, since our object is the comparison, not the models.

\paragraph{Seed-null diagnostic.} For each data set we also compare XGBoost with itself, the two runs differing only in their random seed ($30$ pairs). By symmetry, the expected loss differential of such a comparison is zero, while in each run the stochastic subsampling produces a non-zero observed difference. A procedure that frequently declares one seed ``significantly better'' than the other confuses training noise with a real improvement. Because the data are held fixed, this diagnostic measures the sensitivity of each procedure to training randomness only; it is not a full assessment of size, which Sections~\ref{sec:sim1} and~\ref{sec:sim2} provide.

%%REALRESULTS%%

\section{Discussion}\label{sec:discussion}

\subsection{Are machine learning improvements statistically real?}

Our results give a nuanced answer to the question in the title. With spatially or temporally dependent data, the conventional pipeline (random $K$-fold cross-validation followed by a paired $t$-test, an i.i.d. bootstrap, or even the corrected resampled $t$-test) produces false discoveries at rates between $15\%$ and $40\%$ under moderate to strong dependence in our simulations, instead of the nominal $5\%$. Moreover, random cross-validation targets an interpolation estimand that can understate the gain at new places and times by half, and the case studies show that the ranking of the same learners can change with the prediction task: on hourly weather data, XGBoost appears best under random folds, the learners are indistinguishable when filling two-week gaps, and the random forest is significantly better when forecasting later weeks. Many reported ``significant'' improvements on such data are therefore not statistically supported, and some are answers to a different question from the one of interest.

Switching to blocked cross-validation fixes the estimand but not, by itself, the inference: applied to blocked folds, procedures that resample observations or that model only temporal dependence are \emph{more} anti-conservative than under random folds ($60$--$80\%$ false positives under strong dependence). Cluster procedures (block bootstrap, random-intercept models) remain anti-conservative whenever the dependence range is not small compared with the blocks, and they ignore the correlation of training-induced fold effects. The blocked paired $t$-test, whose behaviour is explained by Proposition~\ref{prop:pairedt}, is the best conventional procedure but still rejects $11$--$18\%$ of true nulls under moderate and strong dependence.

\subsection{What DA-P adds, and what it does not}

\paragraph{Relative to the Hering--Genton test.} The closest existing procedure is the spatial prediction comparison test of Hering and Genton~\cite{hering2011}, which we extended to space--time. When the training-induced variability is small (Experiment~1 and its variants), both procedures are valid, and Hering--Genton is somewhat more powerful because DA-P spends degrees of freedom on a training term that is nearly zero. When the models are refitted learners with substantial fold-to-fold variability (Experiment~2), Hering--Genton rejects $17$--$23\%$ of true nulls under moderate and strong dependence, against $10$--$11\%$ for DA-P; on real data, it declares one seed of XGBoost significantly better than another in $20\%$ of the pedestrian pairs (DA-P: $3\%$), and reports a significant forecasting difference between XGBoost and logistic regression that DA-P attributes to fold-to-fold variability. The training term, whose estimation and limits are characterised by Proposition~\ref{prop:bias}, is thus what distinguishes the comparison of \emph{learning procedures} from the comparison of fixed predictions.

\paragraph{Relative to the kernel variant.} Our initial kernel version, DA-K, over-rejected under moderate and strong dependence ($10$--$15\%$); the pilot study traced this to replicates in which the training variance was truncated at zero, where the downward bias of the fold-centred kernel estimate was no longer compensated. The parametric variance of DA-P does not suffer from kernel down-weighting, and it remained valid under anisotropy, a smoother covariance, heteroscedasticity, non-stationary ranges and irregular sampling, misspecification costing power rather than validity.

\paragraph{What remains.} DA-P is not exactly calibrated in the most difficult setting we studied: learners compared under moderate or strong dependence, where it rejects $10$--$11\%$ of true nulls ($8.7\%$ with the three-way rule). Proposition~\ref{prop:bias} identifies the reason: whichever covariate field happens to be more informative over the whole sample favours one model in every fold, which acts as a fold effect common to all folds, i.e.\ a training correlation $\rho$ larger than assumed, and such a component cannot be estimated from a single cross-validation run~\cite{bengio2004}. The same mechanism explains the over-rejection of all procedures when fold effects are highly correlated in Section~\ref{sec:sim1rho}. The range ratio, which we initially intended to use as a guard, does not identify these cases (Section~\ref{sec:sim1rule}); it is a useful indicator of how much information a comparison can carry, not of its validity.

\subsection{Recommendations}

From these results we derive the following recommendations for comparing two models on spatiotemporal data.
\begin{enumerate}
\item \emph{Define the estimand first.} Choose the validation design to match the intended use of the model: random cross-validation answers an interpolation question at the scale of single observations, blocked cross-validation an interpolation question at the scale of the blocks, and forward chaining a forecasting question.
\item \emph{Never rely on procedures that treat observations as independent} (i.i.d. bootstrap, $z$- or $t$-tests on pooled observations, McNemar's test) or on the paired $t$-test with random folds. Under dependence, the corrected resampled $t$-test is not a remedy either.
\item \emph{Use blocked cross-validation with a test that accounts for both sources of dependence}, such as DA-P, and report its result under $\rho_0=1/K$ and $\rho_0=0.5$; claim an improvement only when it holds under both (the three-way rule), and report the other cases as indeterminate.
\item \emph{Inspect the diagnostics.} Report the estimated ranges, the effective sample size and the range ratio. A small effective sample (a few dozen, as on the Las Rosas field) means that only large differences can be resolved.
\item \emph{Correct for multiplicity} when several models are compared, and treat comparisons on a single data set as exploratory.
\item \emph{Report intervals, not only $p$-values}, on the scale of the loss, as advocated for large data sets in this journal~\cite{caponero2026}: an improvement can be statistically real and practically negligible, or practically important and statistically unresolved.
\end{enumerate}

\subsection{Limitations and future work}

Our simulation designs are stylised: Gaussian fields on a $12\times12$ or $10\times10$ grid (with one irregular variant), a single learner family in Experiment~2, and $K=5$ folds throughout. The working model~\eqref{eq:model} assumes additive fold effects and a stationary loss differential, which may fail when one model is better in some regions and worse in others; a version with a spatially varying $\Delta$ is a natural extension. The correlation $\rho$ of training effects is assumed rather than estimated; repeated blocked cross-validation with different block allocations, or nested cross-validation in the spirit of Bates et al.~\cite{bates2024}, could provide information on it and is, in our view, the most promising route to a fully calibrated test. The exponential variogram does not represent periodic dependence such as daily and weekly cycles, which matter for the pedestrian data. We considered additive losses only; non-additive metrics such as the AUC or the F1-score would require linearisation. Finally, the choice of the parametric variance and of the three-way rule was informed by pilot studies on separate seeds and, for the rule, by the failure of a range-based alternative on the main replicates; both deserve confirmation on further designs and data sets.

\section{Conclusion}\label{sec:conclusion}

Whether a machine learning model is ``better'' than another is a statistical claim, and on spatio-temporal data it depends on two choices that are rarely made explicit: the validation design, which defines what is estimated, and the inferential procedure, which defines how its uncertainty is assessed. We formalised the comparison of two models as inference on the mean of a dependent loss-differential field with correlated fold effects, showed why the paired $t$-test on fold scores fails under random cross-validation, and proposed DA-HAC, a test that combines a variogram-calibrated product-kernel variance, a training-variance term and a Satterthwaite reference distribution, all computed from one blocked cross-validation run.

In controlled simulations and with gradient-boosting learners, conventional procedures declared false improvements up to sixteen times more often than their nominal level, and random cross-validation could halve the estimated gain of a genuinely better model. Blocked cross-validation combined with DA-HAC, or with the paired $t$-test on blocked folds, brought the error rate close to the nominal level under weak dependence and reduced it substantially under stronger dependence, without fully restoring it when the dependence range approaches the size of the study domain, a regime that DA-HAC's diagnostics help to detect. On three public data sets, the conclusion of four of the six comparisons depended on the validation design or on the test: an apparent superiority of XGBoost on hourly weather data vanished under temporal blocking, and two smaller gains were declared significant by independence-based procedures but not by DA-HAC, whose diagnostics showed that one of them could not be resolved with the available data. Large improvements, such as those of gradient boosting over linear baselines, were confirmed by every procedure.

Trustworthy machine learning requires not only accurate and explainable models but also trustworthy evidence that one model is better than another. We hope that the procedure, the diagnostics and the open-source implementation presented here will help practitioners in environmental, agricultural, epidemiological and mobility applications to report improvements that are statistically real.

\backmatter

\bmhead{Acknowledgements}
The author thanks the maintainers of the \texttt{nycflights13}, \texttt{tsibble} and \texttt{agridat} R packages and of the Rdatasets collection for making the data used in this study openly available.

\section*{Declarations}

\begin{itemize}
\item \textbf{Funding.} This research received no specific grant from any funding agency in the public, commercial or not-for-profit sectors.
\item \textbf{Conflict of interest.} The author declares no competing interests.
\item \textbf{Ethics approval and consent to participate.} Not applicable.
\item \textbf{Consent for publication.} Not applicable.
\item \textbf{Data availability.} All data sets are public and are redistributed, with their original licences, in the code repository, together with the raw outputs of all experiments.
\item \textbf{Code availability.} The Python package \texttt{dacomp} and all scripts needed to reproduce the simulations, case studies, tables and figures are available at \url{https://github.com/2zalab/Statistics-for-Data-Science-and-AI} under the MIT licence.
\item \textbf{Author contribution.} I.T. conceived the study, developed the methodology and software, carried out the experiments and wrote the manuscript.
\end{itemize}


\bibliography{references}

\end{document}
